%% =====================================================================
\section{Problem Formulation}\label{sec:problem}
%% =====================================================================
\subsection{Information and Policies}
When the UAV chooses $a(t)$, it knows its labels, confidences and link states up to slot $t$, and the outcome of every earlier transmission as reported by the acknowledgments,
\begin{equation}\label{eq:info}
\I(t)=\big\{\lab(s),c(s),g(s)\big\}_{s\le t}\cup\big\{a(s),\hat X(s+1)\big\}_{s<t},
\end{equation}
together with the ground label of every decoded upload.
It does not know the state $X(t)$ and therefore cannot compute the AoII $\Delta(t)$. A policy $\psi$ is a rule that maps the information $\I(t)$ to a probability distribution on the action set $\A$, from which $a(t)$ is drawn. We allow randomized policies because problems with an average constraint may need them to use the budget exactly. A transmission is pointless when the link is blocked, so we consider only policies with $a(t)=0$ whenever $q_1(g(t))=0$, and we denote their set by $\Psi$.

\subsection{The Optimization Problem}
The response teams need a GS status that is correct as often as possible and, when it is wrong, is corrected quickly. The long-run average AoII measures both: it grows with the fraction of slots in which the status is wrong, and it grows faster when the wrong periods are long, since the AoII counts the slots of each wrong period. Under a policy $\psi$, the average AoII and the average energy are
\begin{align}
\bar\Delta^{\psi}&=\limsup_{T\to\infty}\frac1T\sum_{t=0}^{T-1}\E^{\psi}\big[\Delta(t)\big],\label{eq:avgaoii}\\
\bar E^{\psi}&=\limsup_{T\to\infty}\frac1T\sum_{t=0}^{T-1}\E^{\psi}\Big[E_{a(t)}\big(g(t)\big)\Big],\label{eq:avgE}
\end{align}
where $\E^{\psi}$ is the expectation when the actions follow $\psi$, so that $\bar E^{\psi}$ is the left side of the budget~\eqref{eq:budget} under $\psi$. The UAV must keep the average AoII as low as possible without exceeding the budget:
\begin{equation}\label{eq:P}
\textbf{(P)}\qquad \min_{\psi\in\Psi}\ \bar\Delta^{\psi}\qquad\text{subject to}\qquad \bar E^{\psi}\le\bar E.
\end{equation}

Three features make (P) hard. First, it is partially observed: the cost $\Delta(t)$ depends on $X(t)$, which the UAV never sees, so the policy must act on its own inference about the state. Second, the cost has memory: by~\eqref{eq:aoii}, a wrong GS status adds to the AoII in every slot until it is corrected, so an action affects the cost of all the following slots, not only of the next one. Third, the energy constraint couples all slots: energy spent now is not available later, so the value of an update depends on the updates that the budget will allow in the future. The rest of this section deals with the first feature by replacing the unknown AoII with a probability distribution over its values, which turns (P) into a problem on a state that the UAV knows. Our policy handles the other two features approximately: Section~\ref{sec:structure} values each action by its effect on the AoII of the next slot, and Section~\ref{sec:lyap} enforces the budget online with a virtual queue that prices energy according to what has already been spent.

\subsection{A Belief on the AoII}\label{sec:belief}
The information $\I(t)$ grows without bound. The following lemma shows that a much smaller object carries everything the UAV needs to predict the AoII.

\begin{lemma}[Sufficient statistic]\label{lem:suff}
Given $\hat X(t)$, the AoII determines the state: $X(t)=\hat X(t)$ if $\Delta(t)=0$ and $X(t)=1-\hat X(t)$ if $\Delta(t)>0$. Hence the pair $(\hat X(t),\pi_t)$, where
\begin{equation}\label{eq:belief}
\pi_t(\delta)=\Prob\big\{\Delta(t)=\delta\mid\I(t)\big\},\qquad\delta\in\{0,1,\dots\},
\end{equation}
summarizes $\I(t)$ for predicting the future AoII.
\end{lemma}
\begin{proof}
By~\eqref{eq:aoii}, $\Delta(t)>0$ if and only if $X(t)\neq\hat X(t)$, and $X(t)$ is binary. So $(\hat X(t),\Delta(t))$ determines $X(t)$, and the future of $X$, $\hat X$ and $\Delta$ depends on the past only through $X(t)$, $\hat X(t)$ and $\Delta(t)$.
\end{proof}

The belief $\pi_t$ thus answers both questions at once. Its mass at $\delta=0$ is the probability that the GS is correct, and its mass at $\delta>0$ says how long the GS has probably been wrong. For a given $\delta$, we write $x_\delta$ for the state implied by Lemma~\ref{lem:suff}, i.e., $x_0=\hat X(t)$ and $x_\delta=1-\hat X(t)$ for $\delta>0$.

To store the belief with finite memory, we truncate its support at $\Delta_{\max}$ and keep any mass that would exceed $\Delta_{\max}$ at $\Delta_{\max}$. With $\Delta_{\max}=300$ the belief is a vector of $301$ numbers, about $1.2$\,kB in single precision. The mass at $\Delta_{\max}$ is negligible in all our experiments.

\subsection{Belief Update}
The belief changes for three reasons in each slot: the UAV sees a new on-board label, it may receive a ground label, and the scene evolves. Let $\pi^-_t$ denote the belief at the start of slot $t$, before the new frame is classified. The update has three steps.

\emph{Step 1: correction with the on-board label.} By Bayes' rule, the frame changes the belief in proportion to its likelihood under each value of the AoII:
\begin{equation}\label{eq:bayes}
\pi_t(\delta)\propto\pi^-_t(\delta)\,\Prob\big\{\lab(t),c(t)\mid\Delta(t)=\delta,\I(t-1)\big\}.
\end{equation}
Under Assumption~\ref{as:inf}(ii), the label and the confidence depend on the past only through the current state, which is $x_\delta$ by Lemma~\ref{lem:suff}. By~\eqref{eq:lik}, the likelihood in~\eqref{eq:bayes} is therefore proportional to $c_{x_\delta}(t)/\varpi_{x_\delta}$. All hypotheses $\delta>0$ imply the same state $1-\hat X(t)$, so the likelihood takes only two values, one if the GS is right ($\delta=0$) and one if it is wrong ($\delta>0$). Only their ratio matters. Let $d(t)=\ind{\lab(t)\neq\hat X(t)}$ indicate that the label disagrees with the GS. Their ratio is
\begin{equation}\label{eq:Lambda}
\Lambda_t=\frac{c_{1-\hat X(t)}(t)}{c_{\hat X(t)}(t)}\cdot\frac{\varpi_{\hat X(t)}}{\varpi_{1-\hat X(t)}},
\end{equation}
which says how many times more likely the frame is if the GS is wrong than if it is right. By~\eqref{eq:cx}, the first factor is $c(t)/(1-c(t))$ if the label disagrees with the GS and $(1-c(t))/c(t)$ if it agrees. A label that disagrees with the GS thus has $\Lambda_t$ increasing in its confidence: it is evidence that the GS is wrong, and the more so the more confident it is. A label that agrees with the GS is evidence that the GS is right.

Exact Bayes' rule would multiply the mass at $\delta>0$ by $\Lambda_t$ relative to the mass at $\delta=0$. This treats every frame as an independent observation, which overstates the evidence when errors come in runs. If a frame overlaps half of the previous one and the classifier missed the flood in the previous frame, a second miss is far from independent confirmation (Fig.~\ref{fig:inference}(c)). We therefore temper the ratio with an exponent $\nu\in(0,1]$. With $\mu^-_t=\pi^-_t(\Delta>0)$ the prior probability that the GS is wrong, the corrected belief is
\begin{equation}\label{eq:post}
\pi_t(\delta)=\frac{\pi^-_t(\delta)}{1-\mu^-_t+\mu^-_t\Lambda_t^{\nu}}\times\begin{cases}1, & \delta=0,\\ \Lambda_t^{\nu}, & \delta>0.\end{cases}
\end{equation}
The exponent counts each frame as a fraction $\nu$ of an independent observation, and $\nu=1$ recovers Bayes' rule. A rough guide is the inverse of the mean length of an error run: our on-board errors last $2.2$ slots on average, which suggests $\nu\approx0.45$. In the evaluation we select $\nu\in\{0.25,0.5\}$ on tuning data together with the other parameters of the controller. We also include, as an option, the symmetric update that ignores the prior and uses $\varpi_0=\varpi_1$ with $\nu=1$, in which the confidence itself is the likelihood.

Step 1 multiplies the odds that the GS is wrong by $\Lambda_t^{\nu}$. It moves mass between ``the GS is right'' and ``the GS is wrong'' but does not change the shape of the belief within $\delta>0$.

\emph{Step 2: correction with the ground label.} If the UAV uploaded the frame and the GS decoded it, the acknowledgment carries the ground label, and~\eqref{eq:post} is applied once more with $\nu=1$. Its error probability $\ee_{\rm g}=\ee_{\rm g}(c(t))$ is defined given the state, so no prior correction is needed, and the ratio is $(1-\ee_{\rm g})/\ee_{\rm g}$ if the ground label disagrees with $\hat X(t)$ and $\ee_{\rm g}/(1-\ee_{\rm g})$ otherwise.

\emph{Step 3: prediction.} After the transmission, the GS holds $\xh'=\hat X(t+1)$, which the UAV knows from the acknowledgment or from the absence of one. For each $\delta$, the state $X(t)=x_\delta$ is known, so the UAV knows whether the GS status is now correct for that hypothesis. By~\eqref{eq:aoii}, the AoII grows to $\delta+1$ if $X(t+1)\neq\xh'$ and resets to 0 otherwise. If $x_\delta=\xh'$, the status becomes wrong only if the state changes, which happens with probability $p_{x_\delta}$. If $x_\delta\neq\xh'$, it stays wrong if the state does not change, with probability $1-p_{x_\delta}$. Hence
\begin{equation}\label{eq:pred}
\Prob\big\{\Delta(t+1)=\delta+1\mid\Delta(t)=\delta\big\}=\begin{cases}p_{x_\delta}, & x_\delta=\xh',\\ 1-p_{x_\delta}, & x_\delta\neq\xh',\end{cases}
\end{equation}
with $\Delta(t+1)=0$ otherwise. Summing over $\delta$ under $\pi_t$ gives $\pi^-_{t+1}$. A delivered label that flips the GS status swaps the two cases of~\eqref{eq:pred}, which is how a correction moves the mass of the belief back towards $\delta=0$.

Two summaries of the belief appear in the rest of the paper:
\begin{equation}\label{eq:musum}
\mu_t=\pi_t(\Delta>0),\qquad S_t=\E_{\pi_t}\big[(\Delta+1)\ind{\Delta>0}\big].
\end{equation}
The first is the probability that the GS is wrong. The second weighs each wrong hypothesis by the AoII it will have in the next slot if the GS stays wrong. We write $\mu^-_t$ and $S^-_t$ for the same summaries of $\pi^-_t$.

\subsection{A Worked Example}
\begin{example}\label{ex:belief}
Let the GS report ``not flooded'' ($\hat X(t)=0$), and let the UAV hold the belief $\pi^-_t(0)=0.95$, $\pi^-_t(1)=0.05$: it is fairly sure the GS is right. A frame then arrives with label ``flooded'' and confidence $c=0.9$, so $d(t)=1$. With $\varpi_0=0.81$ and $\varpi_1=0.19$,~\eqref{eq:Lambda} gives $\Lambda_t=(0.81/0.19)(0.9/0.1)=38.4$: the frame is $38$ times more likely if the GS is wrong. With exact Bayes' rule ($\nu=1$),~\eqref{eq:post} gives $\mu_t=0.05\cdot38.4/(0.95+0.05\cdot38.4)=0.67$: one confident frame makes the GS more likely wrong than right. With $\nu=0.5$ the same frame gives only $\mu_t=0.25$, and with $\nu=0.25$ it gives $\mu_t=0.12$. Suppose the UAV stays silent. The prediction~\eqref{eq:pred} with $p_0=0.039$ and $p_1=0.166$ gives $\mu^-_{t+1}=0.23$ for $\nu=0.5$. If the next frame again says ``flooded'' with $c=0.9$, the belief rises to $\mu_{t+1}=0.65$. With the tempered update, the UAV therefore needs two consistent frames to become as convinced as one frame makes it under exact Bayes' rule. Section~\ref{sec:structure} continues this example to show what the UAV does in each of these two slots.
\end{example}

\subsection{Equivalent Belief-State Problem and Solution Approach}
By Lemma~\ref{lem:suff}, (P) is equivalent to a constrained Markov decision process whose state is the belief together with the observable quantities, $(\pi^-_t,\hat X(t),\lab(t),c(t),g(t))$. The per-slot cost is the expected AoII $\E[\Delta(t+1)\mid\I(t),a(t)]$, which the belief gives, and the per-slot energy is $E_{a(t)}(g(t))$. The transition kernel follows from the model: the next label and confidence are drawn from their law given $X(t+1)$; $g(t+1)$ follows the trajectory and the blockage chain~\eqref{eq:block}; $\hat X(t+1)$ follows~\eqref{eq:est}; and $\pi^-_{t+1}$ follows the three steps of Section~\ref{sec:belief}.

The standard approach prices the energy constraint with a Lagrange multiplier $\lambda\ge0$, solves the unconstrained problem with per-slot cost $\E[\Delta(t+1)]+\lambda E_{a(t)}(g(t))$ by value iteration, and searches over $\lambda$ to meet the budget. This is impractical on a UAV: the belief is a point in a high-dimensional simplex, the kernel needs the law of the confidence, the blockage statistics and the trajectory, which change between missions, and the right multiplier depends on all of them. We therefore take a different route. Section~\ref{sec:structure} characterizes the value of each action in one slot in closed form, and Section~\ref{sec:lyap} builds an online policy on this characterization, with an energy price that adapts to the energy actually spent.
